% 16
\begin{frame}{Шаг 1: применяем гладкость}
\small
К точкам $x_k$ и $x_{k+1}$:
\[
\begin{aligned}
f(x_{k+1})
\le{}&
f(x_k)
+\grad f(x_k)^\top(x_{k+1}-x_k)\\
&+\frac L2\norm{x_{k+1}-x_k}^2.
\end{aligned}
\]
Вычитаем $f^\star$:
\[
\gap{x_{k+1}}
\le
\gap{x_k}
+\grad f(x_k)^\top(x_{k+1}-x_k)
+\frac L2\norm{x_{k+1}-x_k}^2.
\]
\end{frame}

% 17
\begin{frame}{Шаг 2: используем выпуклость}
\small
Из слайда 11:
\[
\gap{x_k}
\le
\grad f(x_k)^\top(x_k-x^\star).
\]
Подставляем:
\[
\begin{aligned}
\gap{x_{k+1}}
\le{}&
\grad f(x_k)^\top(x_k-x^\star)\\
&+\grad f(x_k)^\top(x_{k+1}-x_k)\\
&+\frac L2\norm{x_{k+1}-x_k}^2.
\end{aligned}
\]
\end{frame}

% 18
\begin{frame}{Два скалярных произведения объединяются}
\small
Первые два члена имеют общий градиент:
\[
\begin{aligned}
&\grad f(x_k)^\top(x_k-x^\star)
+\grad f(x_k)^\top(x_{k+1}-x_k)\\
&\qquad=
\grad f(x_k)^\top(x_{k+1}-x^\star).
\end{aligned}
\]
Поэтому
\[
\boxed{
\gap{x_{k+1}}
\le
\grad f(x_k)^\top(x_{k+1}-x^\star)
+\frac L2\norm{x_{k+1}-x_k}^2.
}
\]
\end{frame}

% 19
\begin{frame}{Шаг 3: заменяем градиент через шаг}
\small
Используем
\[
\grad f(x_k)=L(x_k-x_{k+1}).
\]
Получаем
\[
\begin{aligned}
\gap{x_{k+1}}
\le{}&
L(x_k-x_{k+1})^\top(x_{k+1}-x^\star)\\
&+\frac L2\norm{x_k-x_{k+1}}^2.
\end{aligned}
\]
Осталось распознать чисто геометрическое тождество.
\end{frame}

% 20
\begin{frame}{Трёхточечное тождество}
\small
Для любых $a,b,c$:
\[
\boxed{
2(a-b)^\top(b-c)
=
\norm{a-c}^2
-\norm{a-b}^2
-\norm{b-c}^2.
}
\]
В нашем случае:
\[
a=x_k,
\qquad
b=x_{k+1},
\qquad
c=x^\star.
\]
\end{frame}

% 21
\begin{frame}{Почему тождество верно}
\small
Заметим:
\[
a-c=(a-b)+(b-c).
\]
Тогда
\[
\begin{aligned}
\norm{a-c}^2
&=\norm{a-b}^2+\norm{b-c}^2\\
&\quad+2(a-b)^\top(b-c).
\end{aligned}
\]
Перенос двух первых слагаемых даёт нужную формулу.

\begin{block}{Никакой новой теории}
Это просто раскрытие квадрата нормы суммы двух векторов.
\end{block}
\end{frame}

% 22
\begin{frame}{Ключевое неравенство для выпуклого случая}
\small
Подставляя трёхточечное тождество, получаем
\[
\begin{aligned}
\gap{x_{k+1}}
\le
\frac L2\Bigl(&
\norm{x_k-x^\star}^2
-\norm{x_{k+1}-x^\star}^2\\
&-\norm{x_k-x_{k+1}}^2\Bigr)
+\frac L2\norm{x_k-x_{k+1}}^2.
\end{aligned}
\]
Последние два члена сокращаются:
\[
\boxed{
\gap{x_{k+1}}
\le
\frac L2
\left(
\norm{x_k-x^\star}^2-
\norm{x_{k+1}-x^\star}^2
\right).
}
\]
\end{frame}

% 23
\begin{frame}{Теперь складываем неравенства}
\small
Для $k=0,\ldots,N-1$:
\[
\sum_{k=0}^{N-1}
\bigl(f(x_{k+1})-f^\star\bigr)
\le
\frac L2
\sum_{k=0}^{N-1}
\left(
\norm{x_k-x^\star}^2-
\norm{x_{k+1}-x^\star}^2
\right).
\]
Справа возникает телескопическая сумма.
\end{frame}

% 24
\begin{frame}{Телескопическая сумма}
\centering
\includegraphics[width=0.88\linewidth]{telescoping_sum.pdf}

\vspace{0.2em}
\small
В результате остаются только начальное и конечное расстояния.
\end{frame}

% 25
\begin{frame}{После суммирования}
\small
Получаем
\[
\sum_{k=1}^{N}
\bigl(f(x_k)-f^\star\bigr)
\le
\frac L2
\left(
\norm{x_0-x^\star}^2-
\norm{x_N-x^\star}^2
\right).
\]
Поскольку норма неотрицательна,
\[
\boxed{
\sum_{k=1}^{N}
\bigl(f(x_k)-f^\star\bigr)
\le
\frac L2
\norm{x_0-x^\star}^2.
}
\]
\end{frame}

% 26
\begin{frame}{Используем монотонность значений функции}
\small
По лемме об убывании
\[
f(x_N)
\le
f(x_{N-1})
\le\cdots\le
f(x_1).
\]
Значит каждое из первых $N$ отклонений не меньше последнего:
\[
 f(x_k)-f^\star
 \ge
 f(x_N)-f^\star.
\]
Поэтому
\[
\boxed{
N\bigl(f(x_N)-f^\star\bigr)
\le
\sum_{k=1}^{N}
\bigl(f(x_k)-f^\star\bigr).
}
\]
\end{frame}

% 27
\begin{frame}{Теорема: скорость GD для выпуклой гладкой функции}
\small
\begin{block}{Теорема}
Пусть $f$ выпукла и $L$-гладка, а
\[
x_{k+1}=x_k-\frac1L\grad f(x_k).
\]
Тогда для любого $N\ge1$
\[
\boxed{
 f(x_N)-f^\star
 \le
 \frac{L\norm{x_0-x^\star}^2}{2N}.
}
\]
\end{block}
Это и есть скорость порядка $O(1/N)$.
\end{frame}

% 28
\begin{frame}{Доказательство теоремы — в трёх строках}
\footnotesize
Из ключевого неравенства:
\[
\sum_{k=1}^{N}
\bigl(f(x_k)-f^\star\bigr)
\le
\frac L2\norm{x_0-x^\star}^2.
\]
Из монотонности:
\[
N\bigl(f(x_N)-f^\star\bigr)
\le
\sum_{k=1}^{N}
\bigl(f(x_k)-f^\star\bigr).
\]
Следовательно,
\[
N\bigl(f(x_N)-f^\star\bigr)
\le
\frac L2\norm{x_0-x^\star}^2,
\]
откуда делением на $N$ получается утверждение.
\end{frame}

% 29
\begin{frame}{Что означает $O(1/k)$?}
\centering
\includegraphics[width=0.70\linewidth]{one_over_k.pdf}

\vspace{0.2em}
\small
При удвоении числа итераций гарантированная верхняя оценка примерно уменьшается вдвое.
\end{frame}

% 30
\begin{frame}{Сколько итераций нужно для точности $\varepsilon$?}
\small
Хотим обеспечить
\[
f(x_k)-f^\star\le\varepsilon.
\]
Из теоремы достаточно потребовать
\[
\frac{L\norm{x_0-x^\star}^2}{2k}
\le
\varepsilon.
\]
Следовательно,
\[
\boxed{
k
\ge
\frac{L\norm{x_0-x^\star}^2}{2\varepsilon}.
}
\]
То есть сложность по точности — порядка $O(1/\varepsilon)$.
\end{frame}
